\chapter{Introduction}
\label{chap:introduction}

\section{Motivation}
\label{sec:motivation}

The proliferation of online lecture content since 2020 has created an
unprecedented body of educational material that is largely inaccessible
by topic.
By 2025, YouTube alone hosted over 800~million videos, with educational content
accounting for one of the fastest-growing categories~\citep{youtube2024stats}.
Massive Open Online Courses (MOOCs) on platforms such as Coursera, edX, and
MIT OpenCourseWare collectively offer tens of thousands of hour-long lecture
recordings spanning every academic discipline.
Yet students navigating this content face a fundamental discoverability problem:
without fine-grained, automatically generated chapter structure, finding the
10-minute explanation of \emph{eigenvalue decomposition} within a 3-hour
linear algebra recording requires either sequential viewing or prior knowledge
of where to look.

Existing video-retrieval systems typically rely on manually authored chapter
markers---available for fewer than 15\,\% of YouTube educational videos at
time of writing~\citep{todo:youtube_chapter_stats}---or on coarse metadata
(title, description) that rarely reflects intra-video structure.
Automated lecture topic segmentation addresses this gap by identifying, from
the video signal itself, the timestamps at which the lecture transitions between
distinct pedagogical units.

\section{Problem statement}
\label{sec:problem}

\emph{Lecture topic segmentation} is the task of partitioning a video
transcript (and its aligned multimodal streams) into a sequence of
contiguous, semantically coherent segments such that each segment covers
a single coherent topic or pedagogical unit.
Formally, given a sequence of $N$ sentences $\{s_1, \ldots, s_N\}$ with
associated timestamps, the task is to predict a set of boundary indices
$\mathcal{B} \subset \{1, \ldots, N-1\}$ that maximises some measure of
intra-segment coherence and inter-segment contrast.

\emph{Hierarchical} lecture topic segmentation extends this to two levels:
\emph{chapters} (coarser, corresponding to major topic shifts) and
\emph{subtopics} (finer, corresponding to transitions between sub-ideas within
a chapter).
The nesting constraint requires that every chapter boundary is also a subtopic
boundary: $\mathcal{B}_{\text{chapter}} \subseteq \mathcal{B}_{\text{subtopic}}$.
This richer representation directly mirrors the structure that expert instructors
use when planning lectures and that viewers expect when navigating course content.

\section{Research questions}
\label{sec:rqs}

\begin{description}
  \item[RQ1] Does a reliability-weighted modality fusion that dynamically
    suppresses uninformative channels outperform fixed-weight fusion and
    text-only baselines?

  \item[RQ2] Does a two-stage boundary predictor (broad candidate generation
    followed by multimodal refinement) achieve better precision--recall balance
    than a single-stage threshold method?

  \item[RQ3] Does explicit hierarchical modelling with a nesting constraint
    improve chapter-level and subtopic-level segmentation quality simultaneously?

  \item[RQ4] Can a small local language model (Llama-3.1-8B) meaningfully
    improve boundary alignment and generate useful segment titles?

  \item[RQ5] Are the gains of the proposed method over the best prior baseline
    statistically significant across the \dataset{} benchmark?
\end{description}

\section{Contributions}
\label{sec:contributions}

This thesis makes the following contributions:

\begin{enumerate}
  \item \textbf{\dataset{}}: A manually annotated benchmark of 30 academic
    lecture videos from five domains, with creator-validated chapter boundaries
    and human-reviewed subtopic annotations, totalling approximately 55~hours
    of transcribed audio.
    (Chapter~\ref{chap:methodology}, Section~\ref{sec:dataset}.)

  \item \textbf{Reliability-Weighted Fusion (N2)}: A modality fusion module
    that assigns each modality a weight inversely proportional to the entropy
    of its gap-score distribution, requiring no labelled data or manual tuning.
    (Chapter~\ref{chap:methodology}, Section~\ref{sec:fusion}.)

  \item \textbf{Two-Stage Boundary Predictor (N1)}: A predictor that first
    generates high-recall candidates using text-only cosine depth scores, then
    re-ranks them with the fused multimodal signal.
    (Chapter~\ref{chap:methodology}, Section~\ref{sec:boundary}.)

  \item \textbf{Hierarchical Segmenter (N3)}: A segmenter that produces
    valid nested chapter and subtopic boundaries by construction, enforcing
    $\mathcal{B}_{\text{chapter}} \subseteq \mathcal{B}_{\text{subtopic}}$
    without joint optimisation.
    (Chapter~\ref{chap:methodology}, Section~\ref{sec:hierarchy}.)

  \item \textbf{Local-LLM Refinement (N4)}: A post-processing module that uses
    an on-device Llama-3.1-8B to shift boundaries by up to $\pm 2$ sentences
    toward cleaner semantic transitions and to generate concise segment titles.
    (Chapter~\ref{chap:methodology}, Section~\ref{sec:refine}.)

  \item \textbf{Comprehensive evaluation}: A comparison of eight methods across
    five metrics with bootstrap confidence intervals and paired statistical
    testing, plus a qualitative error analysis.
    (Chapter~\ref{chap:results}.)

  \item \textbf{Open-source implementation}: All code, the \dataset{} dataset,
    and model weights are released under permissive licences at
    \href{https://huggingface.co/lecseg}{huggingface.co/lecseg} and
    \href{https://zenodo.org/record/lecseg}{zenodo.org/record/lecseg}.
    (Chapter~\ref{chap:deliverables} and Appendix~\ref{app:dataset}.)
\end{enumerate}

\section{Thesis structure}
\label{sec:structure}

\Cref{chap:literature} reviews prior work on text-only, multimodal, and
hierarchical segmentation.
\Cref{chap:methodology} describes the \lecseg{} pipeline, its novel modules,
and the \dataset{} corpus.
\Cref{chap:results} reports quantitative results, ablations, statistical tests,
and qualitative error analysis.
\Cref{chap:conclusion} summarises findings and \cref{chap:future} sketches
directions for future research.
